\chapter{Claim Traceability}
\label{app:claim-traceability}

This appendix records how the thesis-facing claim groups connect to the
audited claim inventory and archived, manifest-bound experiment evidence. It is
a traceability aid rather than an additional source of empirical results. The
controlled inventory contains 46 claims: nine implementation claims, thirteen
method claims, eleven quantitative result claims, four exploratory claims,
and nine required limitations.

\section{Claim Groups and Thesis Locations}

\begin{table}[htbp]
  \centering
  \footnotesize
  \caption{Controlled claim groups and their principal locations in the
  thesis.}
  \label{tab:appendix-claim-groups}
  \begin{tabular}{@{}
    >{\raggedright\arraybackslash}p{0.15\textwidth}
    >{\raggedright\arraybackslash}p{0.34\textwidth}
    >{\raggedright\arraybackslash}p{0.43\textwidth}@{}}
    \toprule
    Claim IDs & Principal thesis location & Primary authority class \\
    \midrule
    \texttt{I-01--I-03}
      & \cref{sec:implementation-scope,sec:implementation-certification}
      & Reference and paper-core output provenance in the implementation
        source and certified-wrapper boundary. \\
    \texttt{I-04--I-05}
      & \cref{sec:implementation-state,sec:implementation-control-flow}
      & Shared Step~1/2/3 runner, partial-order state, and ordinary-list
        sibling backend. \\
    \texttt{I-06--I-08}
      & \cref{sec:implementation-validation,sec:implementation-policies,sec:implementation-certification}
      & Diagnostics, deterministic replay, fixed execution policies, and
        pre-call certification. \\
    \texttt{I-09}
      & \cref{sec:implementation-complexity}
      & Dated ordinary-list scope amendment and missing-theoretical-backend
        boundary. \\
    \texttt{M-01--M-08}
      & \cref{sec:method-scope,sec:method-protocol,sec:method-implementations,sec:method-timing,sec:method-scheduling}
      & Pre-specified formal configuration and the reused timing and scheduling
        control flow. \\
    \texttt{M-09--M-13}
      & \cref{sec:method-aggregation,sec:method-evidence,sec:method-validation,sec:method-environment}
      & Archived manifest-bound summaries, manifest, evidence-contract
        validator, and recorded execution environment. \\
    \texttt{R-01--R-11}
      & \cref{sec:results-correctness,sec:results-runtime-size,sec:results-ratios,sec:results-family,sec:results-variability,sec:results-elapsed,sec:results-replication}
      & Formal-experiment validated evidence and generated analysis tables; the
        pilot experiment is used only for the directional replication claim. \\
    \texttt{E-01--E-04}
      & \cref{sec:results-exploratory}
      & Generated structure/runtime and checked-counter relationship tables,
        interpreted descriptively and non-causally. \\
    \texttt{L-01--L-09}
      & \cref{chap:limitations}
      & Scope amendment, fixed protocol, timing boundaries, sampling design,
        and the formal analysis limitations. \\
    \bottomrule
  \end{tabular}
\end{table}

The detailed wording, evidence locator, and required boundary for every ID are
maintained in
\nolinkurl{docs/thesis/claim_to_evidence_table.md}. The audited copy has the
following SHA-256:

\begin{center}
  \footnotesize\ttfamily
  cc88363a5aa25d55a56b614b02826aa1\allowbreak\newline
  862d7c48c8a1a54bd2c66da17715e2ea
\end{center}

The chapter-level audit is recorded in
\nolinkurl{docs/thesis/final_claim_audit.md}; it verifies coverage of all 46
controlled claims and the cross-chapter interpretation boundaries.

\section{Archived Evidence Bindings}

\begin{table}[htbp]
  \centering
  \footnotesize
  \caption{Archived run manifests used by the thesis analysis.}
  \label{tab:appendix-evidence-bindings}
  \begin{tabular}{@{}
    >{\raggedright\arraybackslash}p{0.20\textwidth}
    >{\raggedright\arraybackslash}p{0.31\textwidth}
    >{\raggedright\arraybackslash}p{0.41\textwidth}@{}}
    \toprule
    Evidence role & Execution identifier & Manifest SHA-256 \\
    \midrule
    Pilot experiment
      & \texttt{week11\_pilot\_v1\_\_run003}
      & \texttt{ef6a9d17df644eab8c3284d04dc356ce}\newline
        \texttt{4b68b4efdfc0da7eecb5a9c7136c5141} \\
    Formal experiment
      & \texttt{week12\_formal\_sorting\_v1\_\_run001}
      & \texttt{f42ab0b77cb7ea91ed66c66bc6947a81}\newline
        \texttt{cd8746b5d2b67d526472260d2b47f850} \\
    \bottomrule
  \end{tabular}
\end{table}

Each manifest binds six archive files: raw measurements, case summaries,
group summaries, case audits, the fixed configuration, and the recorded
runtime environment. The separate formal evidence-contract validator
regenerates cases and checked diagnostics, recomputes scheduling and summaries,
and verifies these hashes before analysis. The two thesis-facing formal
experiment result figures are conversions of the validated SVG analysis
outputs. The five explanatory
figures are generated from committed TikZ sources. Source and generated hashes
for all seven figures are recorded in
\nolinkurl{thesis/figures/README.md}.

\section{Interpretation Boundaries}

The traceability record preserves the following boundaries across chapters:

\begin{itemize}
  \item the paper core recovers output from maintained state, while
    oracle-derived sorted output belongs to the reference pipeline;
  \item the formal experiment evaluates oracle-certified valid-input sorting, not
    recognition;
  \item paper/reference ratios compare different timed pipeline scopes and
    are not like-for-like end-to-end speedups;
  \item pilot and formal absolute runtimes are not pooled;
  \item structure and counter relationships are exploratory and non-causal;
    and
  \item the ordinary-list implementation does not implement the historical
    theoretical backend and carries no linear-time or asymptotic-complexity
    claim.
\end{itemize}
